\chapter{Introducción y objetivos}
\pagenumbering{arabic}
\label{ch:introduccion}

\section{Contexto y motivación}
\label{sec:contexto_motivacion}

La red eléctrica constituye, por su propia naturaleza, un medio de comunicación de alcance prácticamente universal: sus líneas y cables llegan físicamente a cada punto de suministro, por lo que pueden transportar, además de energía, señales de telecomunicación con una cobertura total \cite{5764444}. Las comunicaciones por red eléctrica (\textit{Power Line Communications, \gls{plc}}) aprovechan esta circunstancia para transmitir información sobre la infraestructura existente, sin necesidad de desplegar un medio de comunicación adicional ni de depender de un operador de telecomunicaciones externo, lo que les confiere ventajas evidentes en cobertura, coste y disponibilidad \cite{5764443}. Como contrapartida, el canal eléctrico no fue concebido para transmitir datos: su respuesta en frecuencia, su impedancia y su nivel de ruido varían continuamente en función de las cargas conectadas, y el medio es compartido por todos los equipos de la red, lo que condiciona de forma determinante el diseño de cualquier sistema \gls{plc} \cite{5764443, SHARMA2017704}.

En función de la banda de frecuencias y de la velocidad de transmisión, las tecnologías \gls{plc} suelen clasificarse en cuatro grandes familias, que cubren desde las centenas de hercios hasta los 250 MHz: banda ultraestrecha (\gls{unbplc}), banda estrecha (\textit{Narrow Band \gls{plc}, \gls{nbplc}}), que opera por debajo de 500 kHz, banda cuasiancha y banda ancha, esta última entre 1,8 y 250 MHz \cite{SHARMA2017704}. Las tecnologías de banda ancha maximizan la tasa de datos a costa del alcance y se orientan principalmente al acceso a Internet y a las redes domésticas, mientras que en las aplicaciones de las compañías eléctricas la cobertura resulta más importante que la velocidad, por lo que en este ámbito se prefieren de forma generalizada las tecnologías de banda estrecha \cite{5764444}.

La idea de utilizar el cableado eléctrico como canal de comunicación es casi tan antigua como la propia electrificación: ya en la década de 1920 se documentan sistemas de telemetría sobre líneas de alta tensión \cite{SHARMA2017704}. Sobre la red de distribución, sin embargo, la tecnología ha evolucionado en tres generaciones claramente diferenciadas \cite{5764444}. La primera, el \glslink{controlrizado}{control de rizado} (\textit{ripple control}), surge en la década de 1950 y permite, mediante señales de audiofrecuencia inyectadas en media tensión, una comunicación unidireccional de en torno a 1 bit por segundo destinada a la gestión centralizada de cargas y a la conmutación de tarifas. La segunda generación, desarrollada a partir de las décadas de 1980 y 1990 al amparo de la norma EN 50065-1 de \gls{cenelec}, que reserva la banda de hasta 95 kHz a las compañías \glslink{distribuidora}{distribuidoras} \cite{EN50065_1}, introduce la comunicación bidireccional de baja velocidad y hace posible la lectura automática de contadores (\textit{Automatic Meter Reading, \gls{amr}}). Por último, la tercera generación, basada en la modulación multiportadora (\textit{Orthogonal Frequency Division Multiplexing, \gls{ofdm}}), da lugar a estándares de banda estrecha de alta velocidad como \gls{prime} o \gls{gtresplc}, posteriormente recogidos en las normas \gls{ieee} 1901.2-2013 e \gls{itut} G.9903/G.9904 \cite{5764444, SHARMA2017704}. Cada una de estas generaciones ha multiplicado la tasa de datos de la anterior en aproximadamente dos órdenes de magnitud \cite{5764444}, y esta evolución se describe con mayor detalle en el Capítulo \ref{ch:state_of_the_art}.

El salto cualitativo de la tercera generación coincide en el tiempo con el desarrollo del concepto de \glslink{redinteligente}{red eléctrica inteligente} (\textit{Smart Grid}), en la que la energía y la información circulan de forma bidireccional entre todos los agentes del sistema. Entre las prioridades identificadas para su despliegue, la infraestructura de medición avanzada (\textit{Advanced Metering Infrastructure, \gls{ami}}) ocupa el primer lugar, y la medición inteligente (\textit{smart metering}) se considera el primer hito en el camino hacia la red inteligente \cite{5764443}. La sustitución de los contadores electromecánicos, leídos manualmente in situ, por \glslink{contadorint}{contadores inteligentes} con capacidad de comunicación ha permitido automatizar la lectura, aumentar la granularidad temporal de los datos de consumo y habilitar nuevas funciones de \glslink{telegestion}{telegestión}, como el cambio remoto de tarifas, la detección de fraudes o la supervisión de la calidad de suministro \cite{WorldBank_AMI_Survey}. Dado que la red de baja tensión ya alcanza cada contador, la tecnología \gls{nbplc} se convirtió en la opción natural para la comunicación de último tramo (\textit{``Last Mile''}), y España ha sido uno de los países pioneros en su despliegue masivo: el estándar \gls{prime}, impulsado por Iberdrola desde 2007, sostiene uno de los mayores despliegues de medición inteligente basados en \gls{plc} del mundo \cite{SHARMA2017704}.

En esta arquitectura, el elemento central de cada red de baja tensión es el concentrador de datos (\textit{Data Concentrator Unit, \gls{dcu}}), instalado en el Centro de Transformación (\gls{ct}) de media a baja tensión que alimenta a los usuarios. El \gls{dcu} actúa como nodo maestro de la red \gls{plc}: gestiona la topología de la red y el registro de los contadores, solicita y recoge periódicamente sus lecturas (curvas de carga, cierres de facturación, eventos y alarmas), las valida y almacena temporalmente, y las remite a través de una red de área extensa (\textit{Wide Area Network, \gls{wan}}), habitualmente Ethernet, fibra óptica o red celular, a los sistemas centrales de la \glslink{distribuidora}{distribuidora} (\textit{Distribution System Operator, \gls{dso}}) \cite{5764443, SHARMA2017704}. En sentido inverso, canaliza hacia los contadores las órdenes y configuraciones procedentes de dichos sistemas. Además de estas funciones de comunicación, los concentradores actuales tienden a incorporar funciones de supervisión de la red de baja tensión, como la medida de las magnitudes eléctricas en la salida del transformador mediante un contador de balance, que permite contrastar la energía suministrada con la suma de la facturada y detectar pérdidas o fraudes. Esta tendencia a integrar en un único equipo funciones que tradicionalmente se resolvían con dispositivos independientes (concentración de datos, medida, calidad de suministro y comunicaciones) se traduce en una reducción significativa del coste para la \gls{dso} \cite{SHARMA2017704}.

Sin embargo, la forma en que se construyen hoy en día la mayoría de los \glspl{dcu} comerciales no refleja todavía ese grado de integración. Se trata, por lo general, de equipos marcadamente modulares, compuestos por varias placas de circuito impreso (\gls{pcb}) interconectadas: una tarjeta de procesamiento con Linux embebido, uno o varios módems \gls{plc} (en ocasiones uno por fase) y un módulo independiente de metrología, cada uno con sus propias conexiones, planos de masa y elementos de fijación \cite{6925952, 10488780}. Esta aproximación eleva el coste de la lista de materiales (\textit{Bill Of Materials, \gls{bom}}), el volumen total del equipo y la complejidad del montaje y de las pruebas de producción.
% NOTA: la descripción de la construcción de los DCUs comerciales (y de sus fuentes, párrafo siguiente) se basa en la observación directa de equipos del mercado, no en una referencia bibliográfica. Valorar si se añade alguna referencia (hojas de características, notas de aplicación) o se indica expresamente que se trata de una observación propia.

La fuente de alimentación es, probablemente, el subsistema en el que estas limitaciones resultan más evidentes. En muchos equipos, la tensión de red se adapta mediante transformadores de 50 Hz, en ocasiones trifásicos, seguidos de una rectificación y de una etapa de regulación lineal o de convertidores reductores (\textit{buck}) convencionales de baja tensión, que conmutan a frecuencias elevadas. Esta solución desplaza la conmutación al lado de baja tensión, donde su ruido es más fácil de contener, pero a cambio exige un componente magnético de baja frecuencia muy voluminoso y pesado, que condiciona el tamaño de todo el equipo, además del coste de un transformador de baja frecuencia hecho a medida. La alternativa habitual consiste en emplear módulos conmutados \gls{ac}/\gls{dc} comerciales, compactos y ya certificados, pero cuya frecuencia de conmutación, del orden de la decena a la centena de kilohercios, sitúa la componente fundamental o sus armónicos dentro de las bandas de trabajo del propio \gls{plc} (\gls{cenelec}-A, hasta 95 kHz, y \gls{fcc}, hasta 500 kHz). Dado que la fuente y el módem comparten los mismos terminales de red, el ruido de conmutación llega al receptor \gls{plc} sin atenuación alguna y degrada directamente su relación señal a ruido, por lo que estas fuentes requieren un filtrado de interferencias electromagnéticas (\textit{ElectroMagnetic Interference, \gls{emi}}) muy agresivo, que encarece el equipo y ocupa una parte considerable de su volumen. En ambos casos, además, el diseño de la fuente debe compatibilizarse con un rango de tensión de entrada muy amplio, tanto para una red trifásica, como con la operación con una sola fase presente y con una vida útil de 15 años sin mantenimiento, como se justifica en el Capítulo \ref{ch:state_of_the_art}.

Con este punto de partida, en Nogara Technologies (anteriormente \textit{Advanced Engineering Devices}, AED), empresa en la que se ha desarrollado el presente trabajo, se ha abordado el diseño de un \gls{dcu} de nueva generación concebido desde el principio con un enfoque de máxima integración y eficiencia: un único equipo compacto de carril \gls{din} que reúne en la misma placa base la comunicación \gls{plc} trifásica, la metrología trifásica de supervisión, el núcleo de procesamiento y las comunicaciones Ethernet aisladas, alimentado por una fuente conmutada de alta tensión y alta frecuencia (\textit{High-Voltage High-Frequency, \gls{hv}/\gls{hf}}) diseñada a medida, cuya frecuencia de conmutación se sitúa por encima de 500 kHz para no interferir con las bandas \gls{plc}. Este Trabajo Fin de Máster recoge el diseño, el desarrollo y la validación experimental de dicho equipo.

\section{Objetivos}
\label{sec:objetivos}

El objetivo general de este trabajo es el diseño, desarrollo y validación de cara a su certificación de un \gls{dcu} de nueva generación para \glspl{ct}, con comunicación \gls{nbplc} \gls{prime} y metrología trifásica integradas, alimentado por una fuente de alimentación que conmuta fuera de las bandas \gls{plc} y que opera en un amplio rango de tensiones de entrada. Para alcanzarlo, se plantean los siguientes objetivos específicos:

\begin{itemize}
    \item Revisar el estado del arte de las comunicaciones \gls{plc} aplicadas a la medición inteligente, la arquitectura de los \glspl{ct} y de las redes \gls{prime}, las limitaciones de los \glspl{dcu} comerciales y el marco normativo aplicable, con el fin de identificar los requisitos y las oportunidades de mejora de un equipo de nueva generación.

    \item Diseñar una arquitectura \textit{hardware} que integre en una única placa base las funciones principales del equipo (comunicación \gls{plc} trifásica en las bandas \gls{cenelec}-A y \gls{fcc}, metrología trifásica de supervisión, núcleo de procesamiento y comunicaciones Ethernet con aislamiento de muy alta tensión), con el fin de reducir el coste de fabricación, el volumen total y la complejidad de montaje frente a las soluciones modulares, respetando a la vez los requisitos de aislamiento de un equipo de Clase II instalado en un \gls{ct}.

    \item Diseñar una fuente de alimentación \gls{hv}/\gls{hf} a medida que mantenga en todo momento su frecuencia de conmutación por encima de 500 kHz, fuera de las bandas \gls{cenelec}-A y \gls{fcc}; que opere en un amplio rango de tensiones de entrada, incluso con una única fase presente, y soporte los escenarios anómalos de sobretensión y de fallo de neutro; que proporcione las salidas aisladas que requieren los distintos dominios de masa del equipo; y que no degrade la impedancia que el equipo presenta a la red en las bandas \gls{plc}.

    \item Validar el prototipo mediante una campaña de \glslink{precertificacion}{precertificación} en un laboratorio externo, que abarque las emisiones electromagnéticas, la señalización \gls{plc}, los aislamientos, la seguridad eléctrica y los ensayos eléctricos, de inmunidad, mecánicos y climáticos, como paso previo a la certificación definitiva del equipo.

    \item Valorar y verificar experimentalmente que la elección de la topología de la fuente y de su frecuencia de conmutación contribuye de forma notable a mejorar la calidad de la comunicación \gls{plc}, mediante un estudio de sensibilidad que compare el \gls{dcu} desarrollado con un equipo de referencia y con un equipo alimentado por una fuente comercial que conmuta dentro de las bandas \gls{plc}.
\end{itemize}

\section{Estructura de la memoria}
\label{sec:estructura_memoria}

La memoria se organiza en cinco capítulos, además del presente. El Capítulo \ref{ch:state_of_the_art} revisa el estado del arte: la evolución histórica de las comunicaciones \gls{plc}, la arquitectura de los \glspl{ct} y de las redes \gls{prime}, las limitaciones de los equipos comerciales y el marco normativo aplicable. El Capítulo \ref{ch:hardware_design} recoge los requisitos del equipo y describe su arquitectura \textit{hardware}: el esquema de aislamientos galvánicos, el subsistema de procesamiento, el frente analógico y el acoplamiento \gls{plc} trifásico, y el subsistema de metrología y su calibración. El Capítulo \ref{HV_HF_PSU} se dedica a la fuente de alimentación \gls{hv}/\gls{hf}, desde sus requisitos hasta la topología adoptada y su validación preliminar. El Capítulo \ref{ch:precompliance_results} presenta el estudio de sensibilidad de la comunicación \gls{plc} y los resultados de la campaña de \glslink{precertificacion}{precertificación}, y el último capítulo recoge las conclusiones y las líneas de trabajo futuro. Finalmente, el Anexo \ref{anx:ensayos_complementarios} reúne los ensayos complementarios de la campaña cuyo resultado se limita a la verificación de un criterio de funcionamiento.